% ==========================================================================
% GENERATED by scripts/split_paper.py. Do not edit this file.
%
% Source: imported_papers/grokking/iclr2027_conference_v2.tex
% Part 11: Neuron deletion and adaptation
%
% Edits belong upstream, in the published paper. To change what the
% chapter includes, change imports/grokking.toml and re-run
% the splitter.
% ==========================================================================
\section{Neuron deletion and adaptation}
\label{app:consequences}

\subsection{Neuron deletion}
\label{app:deletion}
Deleting a first-hidden-layer neuron sets its outgoing weights to zero; the network is then evaluated on the test set without retraining. Random deletion removes fractions $0, 0.1, \ldots, 0.8$ of the layer, with $20$ random masks per fraction ($30$ in the ten-seed experiment of Figure~\ref{fig:main}c). Deletion tolerance is the area under the curve of test accuracy against the deleted fraction at these nine fractions, computed by the trapezoidal rule and divided by the range $0.8$, so that it lies between $0$ and $1$. Norm-ranked deletion removes neurons in decreasing order of outgoing norm at the same fractions. Greedy deletion evaluates every remaining neuron, removes the one whose removal lowers test accuracy the most, and repeats until accuracy falls below $0.5$; we report the number of neurons removed. Greedy deletion selects neurons by their effect on the test set rather than by their weights.

Table~\ref{tab:targeted} compares the reference, full credit equalization and DFC under random and norm-ranked deletion in one matched three-seed experiment (Figure~\ref{fig:fragile}). Equalization improves tolerance under both procedures, and the improvement is larger under norm-ranked deletion, which removes exactly the neurons a concentrated circuit relies on. DFC lies between the two. Table~\ref{tab:plane} reports greedy deletion for the reference and three equalization strengths: the number of neurons that must be removed grows with participation, from $23$ at $\beta = 0$ to $61$ at $\beta = 1$. The deletion advantage of the distributed circuit thus does not depend on selecting neurons by their norm. Over ten seeds, random-deletion tolerance is $0.672\,(0.012)$ for the reference and $0.773\,(0.014)$ under full equalization (Figure~\ref{fig:main}c).

\paragraph{Removal in retention order.} Pruning removes neurons in the order that keeps performance. We remove first-hidden-layer neurons in increasing order of outgoing norm (magnitude pruning) and, greedily, the neuron whose removal raises the training cross-entropy least, and record test accuracy without retraining (Table~\ref{tab:retain}; five seeds, the networks of Table~\ref{tab:decay}). The concentrated circuit keeps test accuracy above $0.9$ until $62\%$ of its neurons are removed by magnitude, against $20\%$ for the distributed circuit. Greedy removal takes away $191$ against $160$ neurons before test accuracy falls below $0.9$, and $177$ against $117$ before it falls below $0.95$. The nearly unused neurons of the concentrated circuit can be removed at no cost, so it is the easier circuit to compress.

\begin{table}[!ht]
\caption{Removal in retention order (modular addition, $p=31$, five seeds; standard deviations in parentheses). Magnitude pruning removes neurons in increasing order of outgoing norm, evaluated at fractions $0, 0.1, \ldots, 0.8$; greedy pruning removes the neuron whose removal raises the training cross-entropy least.}
\label{tab:retain}
\begin{center}
\small
\begin{tabular}{lcc}
\toprule
 & Concentrated ($\beta=0$) & Distributed ($\beta=1$) \\
\midrule
Magnitude pruning, tolerance & $0.949$ & $0.718$ \\
Magnitude pruning, largest fraction with test $\geq 0.9$ & $0.62\,(0.04)$ & $0.20\,(0.00)$ \\
Greedy pruning, neurons removed with test $\geq 0.9$ & $191\,(4)$ & $160\,(18)$ \\
Greedy pruning, neurons removed with test $\geq 0.95$ & $177\,(5)$ & $117\,(37)$ \\
\bottomrule
\end{tabular}
\end{center}
\end{table}

\begin{figure}[!ht]
\begin{center}\includegraphics[width=0.9\linewidth]{figures/fig4_fragility.pdf}\end{center}
\caption{Neuron deletion for the reference, full credit equalization ($\beta=1$) and DFC on modular addition ($p=31$, three seeds). \textbf{(a)} Test accuracy against the fraction of randomly deleted first-hidden-layer neurons. \textbf{(b)} Deletion tolerance, the normalized area under the accuracy curve over the deletion fractions, under random and norm-ranked deletion.}
\label{fig:fragile}
\end{figure}

\begin{table}[!ht]
\caption{Deletion tolerance under random and norm-ranked deletion (modular addition, $p=31$, three seeds). Norm-ranked deletion removes neurons in decreasing order of outgoing norm.}
\label{tab:targeted}
\begin{center}
\small
\begin{tabular}{lccc}
\toprule
Deletion procedure & Reference & Credit equalization ($\beta{=}1$) & DFC \\
\midrule
Random        & $0.674$ & $\mathbf{0.770}$ & $0.725$ \\
Norm-ranked   & $0.332$ & $\mathbf{0.581}$ & $0.518$ \\
\bottomrule
\end{tabular}
\end{center}
\end{table}

\begin{table}[!ht]
\caption{Greedy neuron deletion on modular addition ($p=31$, three seeds). At each step the procedure removes the neuron whose removal lowers test accuracy the most. The last column gives the mean number of neurons removed before test accuracy falls below $0.5$. Standard deviations of circuit participation are in parentheses.}
\label{tab:plane}
\begin{center}
\small
\begin{tabular}{lccc}
\toprule
Condition & Test & Circuit PR & Neurons removed \\
\midrule
Reference & $0.998$ & $0.595\,(0.023)$ & $23.3$ \\
Credit equalization, $\beta{=}0.25$ & $0.997$ & $0.940\,(0.017)$ & $43.0$ \\
Credit equalization, $\beta{=}0.5$ & $0.986$ & $0.975\,(0.002)$ & $59.0$ \\
Credit equalization, $\beta{=}1$ & $0.976$ & $0.992\,(0.002)$ & $60.7$ \\
\bottomrule
\end{tabular}
\end{center}
\end{table}

\subsection{Minimal circuits}
\label{app:minimal}
We prune the first hidden layer of each trained multilayer perceptron, or the MLP of the transformer, and record the test-accuracy curve along the pruning order and the number of neurons kept just before it first falls below a threshold. Greedy pruning removes at each step the neuron whose removal raises the training cross-entropy least, recomputed after every removal, so the order uses only training pairs and the held-out pairs only set where the curve is read. Magnitude pruning removes neurons in increasing order of outgoing norm. All networks are trained on modular addition at $p = 31$ with five fresh seeds. The network with a decayed per-neuron gain multiplies each first-layer pre-activation by a trainable gain with its own bias, both decayed like the weights, and the transformer is the wide-MLP configuration of Appendix~\ref{app:transformer}. Greedy pruning to test accuracy $0.95$ changes the mean trigonometric fit of Appendix~\ref{app:algo} by at most $0.11$ in the multilayer perceptrons, for example from $0.66$ to $0.55$ for the reference and from $0.90$ to $0.81$ under equalization at width $2048$, so their pruned circuits keep the Fourier structure of the full networks (Table~\ref{tab:minimalcheck}). In the transformer the fit falls from $0.99$ to $0.44$ in both conditions, which lies in the range $0.42$ to $0.68$ of the full networks of the grokking rules in Appendix~\ref{app:algo}.

\begin{table}[!ht]
\caption{Number of neurons that pruning keeps at test accuracy $0.90$ or $0.95$ (modular addition, $p=31$, five seeds; range over seeds in brackets).}
\label{tab:minimal}
\begin{center}
\footnotesize
\setlength{\tabcolsep}{4pt}
\begin{tabular}{llccccc}
\toprule
Network & Width & Circuit PR & Greedy, $0.90$ & Greedy, $0.95$ & Magnitude, $0.95$ \\
\midrule
Adam with L2, reference & $256$ & $0.594$ & $59$ & $70$ $[63, 81]$ & $88$ \\
Adam with L2, equalized & $256$ & $0.993$ & $96$ & $155$ $[100, 230]$ & $215$ \\
Decayed per-neuron gain & $256$ & $0.096$ & $23$ & $24$ $[22, 25]$ & $25$ \\
AdamW, reference & $256$ & $0.931$ & $61$ & $77$ $[67, 96]$ & $142$ \\
AdamW, equalized & $256$ & $0.970$ & $61$ & $115$ $[71, 171]$ & $217$ \\
Adam-mini with L2 & $256$ & $0.586$ & $61$ & $72$ $[55, 90]$ & $89$ \\
\midrule
Adam with L2, reference & $2048$ & $0.168$ & $58$ & $73$ $[33, 98]$ & $143$ \\
Adam with L2, equalized & $2048$ & $0.999$ & $132$ & $171$ $[138, 229]$ & $1692$ \\
Decayed per-neuron gain & $2048$ & $0.009$ & $16$ & $18$ $[14, 22]$ & $18$ \\
\midrule
Transformer, reference & $4096$ & $0.609$ & $111$ & $140$ $[89, 171]$ & $258$ \\
Transformer, equalized & $4096$ & $0.994$ & $395$ & $462$ $[416, 536]$ & $1120$ \\
\bottomrule
\end{tabular}
\end{center}
\end{table}

At width $2048$ the greedy circuit at threshold $0.95$ under equalization is $7.1$ to $13.5$ times that of the gain network. The order gain network, reference, equalization holds in every seed at thresholds $0.90$ and $0.95$ and in both pruning orders. With the architecture fixed, equalization enlarges the greedy circuit $2.4$ times and the magnitude circuit $12$ times. In the distributed networks greedy pruning finds far smaller circuits than magnitude pruning, because nearly equal outgoing norms do not indicate which neurons the computation needs. Participation and minimal circuit size measure different things: the AdamW reference has nearly equal outgoing norms but prunes to about as few neurons as the Adam reference.

\paragraph{Starting accuracy and stopping point.} The networks reach different test accuracies before pruning, and a network that starts closer to the threshold crosses it after fewer removals. We therefore also stop one or two points below each network's own test accuracy (Table~\ref{tab:minimalcheck}). At width $2048$ the reference and equalization both start at $0.998$, and the order gain network, reference, equalization holds in every seed at both drops and in both pruning orders; two points below, the greedy circuit under equalization is $2.6$ times that of the reference. At width $256$, where equalization starts lower, it is $1.8$ times that of the reference and larger in every seed. The gain network starts lower than the reference, which would enlarge its circuit rather than shrink it. Taking the smallest size at which the curve is still at or above $0.95$, instead of the first crossing, changes the greedy means at width $2048$ from $73$, $171$ and $18$ to $71$, $162$ and $17$.

\begin{table}[!ht]
\caption{Pruning relative to each network's own test accuracy (same networks as Table~\ref{tab:minimal}, means over five seeds). Test is the accuracy before pruning; the next two columns give the neurons that greedy pruning keeps one and two points below it. Fit is the trigonometric fit of Appendix~\ref{app:algo} for the full network and for the greedy circuit at test accuracy $0.95$.}
\label{tab:minimalcheck}
\begin{center}
\footnotesize
\setlength{\tabcolsep}{4pt}
\begin{tabular}{llccccc}
\toprule
Network & Width & Test & Greedy, $-0.01$ & Greedy, $-0.02$ & Fit, full & Fit, pruned \\
\midrule
Adam with L2, reference & $256$ & $0.998$ & $99$ & $82$ & $0.54$ & $0.47$ \\
Adam with L2, equalized & $256$ & $0.980$ & $199$ & $146$ & $0.47$ & $0.44$ \\
Decayed per-neuron gain & $256$ & $0.975$ & $25$ & $25$ & $0.49$ & $0.48$ \\
AdamW, reference & $256$ & $0.992$ & $109$ & $95$ & $0.36$ & $0.42$ \\
AdamW, equalized & $256$ & $0.967$ & $116$ & $88$ & $0.43$ & $0.41$ \\
Adam-mini with L2 & $256$ & $0.995$ & $92$ & $83$ & $0.58$ & $0.51$ \\
\midrule
Adam with L2, reference & $2048$ & $0.998$ & $128$ & $97$ & $0.66$ & $0.55$ \\
Adam with L2, equalized & $2048$ & $0.998$ & $342$ & $252$ & $0.90$ & $0.81$ \\
Decayed per-neuron gain & $2048$ & $0.972$ & $18$ & $18$ & $0.75$ & $0.75$ \\
\midrule
Transformer, reference & $4096$ & $1.000$ & $364$ & $249$ & $0.99$ & $0.44$ \\
Transformer, equalized & $4096$ & $1.000$ & $1208$ & $867$ & $0.99$ & $0.44$ \\
\bottomrule
\end{tabular}
\end{center}
\end{table}

\subsection{Matching participation through outgoing-weight constraints}
\label{app:suff}
Does the deletion advantage depend on how the distributed circuit was trained, or only on the resulting distribution of outgoing norms? We compare two ways of reaching the same circuit participation. Credit equalization changes the gradients. An \emph{outgoing-norm constraint} leaves the gradients unchanged and instead rescales the outgoing weights of the first hidden layer after every update to a prescribed norm profile with participation $0.60$, $0.80$ or $0.99$; the profile has two norm levels, each shared by half of the neurons. All conditions use five seeds (Table~\ref{tab:suff}).

\paragraph{Deletion.} At participation $0.99$, credit equalization and the outgoing-norm constraint give the same random-deletion tolerance, $0.777$; at participation $0.60$, the constrained network and the reference differ by $0.02$. Greedy deletion removes $66.7\,(11.1)$ neurons under equalization and $68.7\,(10.8)$ under the constraint before accuracy falls below $0.5$. Deletion tolerance thus follows the distribution of outgoing norms in both ways of producing it.

\paragraph{Adaptation.} We also fine-tune these networks from addition to subtraction and measure the epochs to test accuracy $0.9$, evaluated every $20$ epochs. Within the constrained networks, adaptation takes $740$, $1228$ and $1480$ epochs at participation $0.60$, $0.80$ and $0.99$; the equalized network takes $1608$ epochs. Adaptation slows with participation under both interventions. At participation $0.60$, the reference takes $464$ epochs and the constrained network $740$: the reference has $87$ neurons with an outgoing norm below one tenth of the layer mean, which the new task can recruit, whereas the constraint holds every neuron at one of two norms and leaves none. The difference points to the low-norm neurons, which a related task can recruit.

\begin{table}[!ht]
\caption{Credit equalization compared with an outgoing-norm constraint of matched participation (modular addition, $p=31$, five seeds; standard deviations in parentheses). Random del is the random-deletion tolerance; Adapt is the number of fine-tuning epochs to test accuracy $0.9$ on subtraction.}
\label{tab:suff}
\begin{center}
\small
\begin{tabular}{lcccc}
\toprule
Condition & Test & Circuit PR & Random del & Adapt \\
\midrule
Reference            & $0.995\,(0.007)$ & $0.597\,(0.020)$ & $0.669$ & $464$ \\
Credit equalization ($\beta{=}1$)   & $0.976\,(0.015)$ & $0.993\,(0.002)$ & $0.777$ & $1608$ \\
\midrule
Outgoing-norm constraint, $0.60$   & $0.987\,(0.016)$ & $0.600\,(0.000)$ & $0.647$ & $740$ \\
Outgoing-norm constraint, $0.80$   & $0.977\,(0.008)$ & $0.800\,(0.000)$ & $0.675$ & $1228$ \\
Outgoing-norm constraint, $0.99$   & $0.976\,(0.012)$ & $0.990\,(0.000)$ & $0.777$ & $1480$ \\
\bottomrule
\end{tabular}
\end{center}
\end{table}

\subsection{Adaptation to a new operation}
\label{app:adapt}
Networks first grok addition at $p = 31$ at each $\beta$. We then fine-tune the whole network on the new operation with the reference rule, backpropagation with Adam at learning rate $10^{-3}$ and a fresh optimizer state, in every condition, on the new operation's own split with $40\%$ of the pairs held out, and record the number of epochs until test accuracy on the new operation reaches $0.9$, evaluated every $100$ epochs (Table~\ref{tab:adapt}).

\paragraph{Transfer to subtraction.} Adaptation time increases with the participation of the source circuit at every step, from $640$ epochs at $\beta = 0$ to $1540$ at $\beta = 1$; training from scratch takes $1410$ epochs. The seed ranges at the two endpoints do not overlap: $400$ to $800$ epochs at $\beta = 0$ against $1200$ to $1800$ at $\beta = 1$. An independent ten-seed run, evaluated every $20$ epochs, gives $690\,(164)$ and $1550\,(191)$ epochs. Pretraining on addition helps from the concentrated circuit and not from the distributed one.

\paragraph{The role of low-norm neurons.} We count the neurons whose outgoing norm is below one tenth of the layer mean. The reference has about $85$ such neurons. Our hypothesis is that subtraction can recruit them with little disturbance to the addition circuit. Already at $\beta = 0.25$ the count is zero, and adaptation time rises from $640$ to $940$ epochs. It keeps rising to $1540$ epochs at $\beta = 1$ as the circuit spreads further, which is consistent with the new operation having to modify a larger share of a distributed representation. In continual learning, inactive units signal a loss of plasticity that resetting them restores \citep{sokar2023dormant,dohare2024loss}; here the low-norm neurons of the reference circuit coincide with faster adaptation, which suggests that unused capacity can serve a related task when the rest of the circuit stays intact.

\paragraph{Transfer after switching the learning rule.} We train reference networks on addition for $20{,}000$ epochs, then copy each into two branches that continue for $30{,}000$ epochs with fresh Adam states, using either the reference rule or full equalization. Both branches are then fine-tuned identically on subtraction, with evaluation every $20$ epochs. Across five seeds, equalized branches reach $0.9$ test accuracy in $1052\,(273)$ epochs versus $608\,(160)$ for reference continuations; the mean paired difference is $444\,(131)$ epochs. Source accuracies are $0.987$ and $0.996$, respectively. Equalized versus reference continuations have outgoing-weight participation of $0.99$ versus $0.60$ and ablation participation of $0.72$ versus $0.40$. The participation measures show little further change between $10{,}000$ and $30{,}000$ continuation epochs. Parentheses give standard deviations.

\paragraph{Fine-tuning with the source rule.} The transfer runs above fine-tune with the reference rule. To check that the ordering does not depend on the fine-tuning rule, we fine-tune each source network with both rules, evaluated every $20$ epochs (Table~\ref{tab:ownrule}). Over five seeds, the concentrated network reaches $0.9$ on subtraction after $580\,(151)$ epochs with the reference rule and $1000\,(196)$ with full equalization; the distributed network takes $1496\,(195)$ and $1856\,(95)$. The distributed source is slower under both fine-tuning rules and in every seed, by a factor of $2.6$ and $1.9$. The slower adaptation of the distributed circuit does not depend on the fine-tuning rule.

\begin{table}[!ht]
\caption{Fine-tuning from addition to subtraction with the reference rule or with full credit equalization, for a concentrated and a distributed source network (modular addition, $p=31$, five seeds; standard deviations in parentheses). Adaptation time is the number of epochs to test accuracy $0.9$, evaluated every $20$ epochs.}
\label{tab:ownrule}
\begin{center}
\small
\begin{tabular}{llccc}
\toprule
Source network & Fine-tuning rule & Source test & Epochs to $0.9$ & Circuit PR after \\
\midrule
Concentrated ($\beta{=}0$) & reference & $0.992$ & $580\,(151)$ & $0.62$ \\
Concentrated ($\beta{=}0$) & equalization & $0.992$ & $1000\,(196)$ & $0.99$ \\
Distributed ($\beta{=}1$) & reference & $0.969$ & $1496\,(195)$ & $0.60$ \\
Distributed ($\beta{=}1$) & equalization & $0.969$ & $1856\,(95)$ & $0.99$ \\
\bottomrule
\end{tabular}
\end{center}
\end{table}

\paragraph{Nearly unused neurons and the kind of target.} We pretrain addition networks at $\beta = 0$ and $\beta = 1$ for $20{,}000$ epochs and fine-tune them as above on five targets, evaluated every $20$ epochs (Table~\ref{tab:spare}, five seeds). A third source is the concentrated network with its nearly unused neurons removed. Their incoming weights, bias and outgoing weights are set to zero, which removes $75$ to $96$ neurons and leaves test accuracy on addition at $0.996$; they stay at zero during fine-tuning, since a ReLU unit with zero incoming weights and bias receives no gradient. Removing them slows adaptation by $552\,(217)$ epochs on subtraction and $492\,(260)$ epochs on $(b - a) \bmod p$, in every seed, which is $67\%$ and $54\%$ of the difference to the distributed network. Two targets keep the hidden features of addition and change only the readout, $(a + b + 7) \bmod p$ and $\pi((a + b) \bmod p)$ for a fixed random permutation $\pi$. Both are learned from every source $2.3$ to $6.6$ times faster than from scratch, and the concentrated network has no advantage.

\begin{table}[!ht]
\caption{Adaptation from addition to five targets, in epochs to test accuracy $0.9$ (modular addition source, $p=31$, five seeds; standard deviations in parentheses). The third column is the concentrated network with its nearly unused neurons removed.}
\label{tab:spare}
\begin{center}
\small
\setlength{\tabcolsep}{4pt}
\begin{tabular}{lcccc}
\toprule
Target & Concentrated & Nearly unused removed & Distributed & From scratch \\
\midrule
$(a - b) \bmod p$ & $652\,(108)$ & $1204\,(303)$ & $1480\,(382)$ & $1300\,(83)$ \\
$(b - a) \bmod p$ & $528\,(72)$ & $1020\,(296)$ & $1440\,(303)$ & $1304\,(46)$ \\
$(a + b + 7) \bmod p$ & $204\,(60)$ & $204\,(69)$ & $224\,(39)$ & $1356\,(97)$ \\
$\pi((a + b) \bmod p)$ & $560\,(115)$ & $452\,(95)$ & $464\,(83)$ & $1308\,(111)$ \\
$(a \cdot b) \bmod p$ & $1940\,(124)$ & $2324\,(224)$ & $1988\,(202)$ & $1148\,(52)$ \\
\bottomrule
\end{tabular}
\end{center}
\end{table}

\paragraph{Transfer follows the circuit.} We also fine-tune the four networks of Table~\ref{tab:decay} on subtraction and on $(b - a) \bmod p$, five seeds each. The two networks with distributed circuits, produced by credit equalization or by per-neuron decay alone, take $1572\,(189)$ and $1408\,(141)$ epochs on subtraction and $1560\,(279)$ and $1452\,(74)$ on $(b - a) \bmod p$. The two networks with concentrated circuits, the reference and equalization with the reference's ratio of gradient to decay, take $616\,(117)$ and $536\,(86)$ on subtraction and $528\,(77)$ and $512\,(92)$ on $(b - a) \bmod p$. On both targets every run from a distributed circuit is slower than every run from a concentrated one. Across these four training conditions, adaptation time thus follows the circuit rather than the rule that produced it.

\paragraph{Transfer to multiplication.} The same pretrained networks take $1800$ to $2020$ epochs to reach the threshold on multiplication, against $1200$ from scratch. Multiplication does not share the operand frequencies of addition, and pretraining slows adaptation at every $\beta$. In our experiments the advantage of the concentrated circuit appears only for the target task that shares the operand frequencies of the source task.

\paragraph{Fine-tuning learning rate.} At one third and at three times the reference learning rate, adaptation time still increases with $\beta$, and the fully equalized start takes $1.9$ to $2.5$ times as long as the concentrated start (Table~\ref{tab:ftlr}). The reverse transfer, from subtraction to addition, shows the same ordering: $740$ epochs at $\beta = 0$ and $1600$ at $\beta = 1$.

\begin{table}[!ht]
\caption{Fine-tuning from addition to a new operation, with standard deviations in parentheses. Subtraction results for $\beta = 0, 0.25, 1$ and from scratch use ten seeds; the intermediate settings and all multiplication results use five. Adaptation time is the number of epochs to test accuracy $0.9$. Low-norm neurons have an outgoing norm below one tenth of the layer mean.}
\label{tab:adapt}
\begin{center}
\footnotesize
\begin{tabular}{lccccc}
\toprule
Condition & Circuit PR & Low-norm neurons & Source test & \multicolumn{2}{c}{Epochs to $0.9$ on} \\
 & & & & subtraction & multiplication \\
\midrule
$\beta=0$ (reference) & $0.593\,(0.034)$ & $85$ & $0.995$ & $640\,(128)$  & $1800\,(141)$ \\
$\beta=0.25$         & $0.933\,(0.020)$ & $0$  & $0.989$ & $940\,(191)$  & $1940\,(233)$ \\
$\beta=0.5$          & $0.972\,(0.006)$ & $0$  & $0.990$ & $1200\,(89)$  & $1980\,(248)$ \\
$\beta=0.75$         & $0.983\,(0.003)$ & $0$  & $0.981$ & $1300\,(141)$ & $2020\,(133)$ \\
$\beta=1$            & $0.991\,(0.002)$ & $0$  & $0.972$ & $1540\,(180)$ & $1940\,(136)$ \\
From scratch         & ---              & ---  & ---     & $1410\,(137)$ & $1200\,(63)$ \\
\bottomrule
\end{tabular}
\end{center}
\end{table}

\begin{table}[!ht]
\caption{Adaptation time from addition to subtraction across an order of magnitude in the fine-tuning learning rate, in epochs to test accuracy $0.9$ (standard deviations in parentheses). The outer columns use three seeds; the middle column repeats Table~\ref{tab:adapt} (ten seeds for the endpoints, five for $\beta = 0.5$).}
\label{tab:ftlr}
\begin{center}
\small
\begin{tabular}{lccc}
\toprule
Start & $3\times 10^{-4}$ & $10^{-3}$ & $3\times 10^{-3}$ \\
\midrule
$\beta=0$ (reference) & $1933\,(525)$ & $640\,(128)$  & $267\,(47)$ \\
$\beta=0.5$          & $3067\,(205)$ & $1200\,(89)$  & $533\,(47)$ \\
$\beta=1$            & $3733\,(262)$ & $1540\,(180)$ & $667\,(94)$ \\
\midrule
ratio $\beta{=}1$ to $\beta{=}0$ & $1.9$ & $2.4$ & $2.5$ \\
\bottomrule
\end{tabular}
\end{center}
\end{table}
